\section{Conjuntos de datos}
\label{sec:met_datasets}

\subsection{HAM10000}
\label{sec:met_ham}

\pendiente{Se presentan la composición del conjunto, la distribución de clases y su desbalance.}

\subsection{Partición disjunta por lesión}
\label{sec:met_particiones}
\pendiente{HAM10000 contiene varias imágenes de una misma lesión y la partición original se estratificó por
imagen, de modo que imágenes de una misma lesión quedaron repartidas entre entrenamiento y test. Esta
subsección explica ese problema, conocido como fuga de información entre particiones (\emph{data leakage}),
lo cuantifica (el 36\% del test comparte lesión con entrenamiento; el 63\% en melanoma) y describe cómo se
midió su efecto sobre las métricas: reevaluación sobre el subconjunto de test cuya lesión no aparece en
entrenamiento ni en validación, y control con un modelo que nunca vio HAM10000 (DermLIP \emph{zero-shot}).

Se describe luego la partición disjunta por lesión que se construyó como reemplazo, con test idéntico al
original y validación agrupada (\texttt{StratifiedGroupKFold}), sus tamaños finales y qué resultados se
recalcularon sobre ella.}

\subsection{SkinCaRe, MM-Skin y HIBA}
\pendiente{Se describe el rol de cada conjunto: entrenamiento de la generación (SkinCAP y SkinCoT),
evaluación con referencia (MM-Skin) y evaluación fuera de distribución (HIBA, solo test). Se indica la
modalidad de cada uno, en particular que SkinCaRe es de fotografía clínica y HAM10000 dermatoscópico, con
remisión a la Sección~\ref{sec:marco_generalizacion}.}
